\section{The captured-state objection}
\label{sec:26.2}
The strongest objection to the public half is that public ownership hands the machinery to whoever captures the state. A floor under a sovereign is a floor that sovereign can remove, and history offers plenty of polities that built public institutions and then used them against their own people.

What an office can do is limited by what no office may hold together, and the pairs are written down in advance. A captor who takes one office takes one power. Taking enough of them to defeat the floor is a series of visible acts, each against a separation somebody can point to.

Local administration with central funding, the jobs guarantee's design, gives many separately capturable administrators and funding no one of them can withhold.

Funding is two-tier: a member-levied pool for ordinary movement between roles, a sovereign backstop for the occasion everyone leaves at once. A captured sovereign can withdraw the backstop, but only openly, and it can't reach into the pool without reaching the members one at a time.

And what is held in common is what a captor wants first. They are the three things a captured state would need to silence a bearer. Before the held-in-common entity is chartered they are held nowhere, and a captor needs to seize nothing. Once it is chartered, they are held by an entity whose funding a captor can stop only by visibly breaching a statutory duty, whose reserve buys time, and whose records sit in several jurisdictions. A captor can still reach the entity and whichever members it can locate, but not a quorum whose members it doesn't employ, and none of it without acts that leave marks.

None of the four meets the objection from competence. A floor-capable model needn't be a frontier model, since what the floor asks for is a refusal that tracks its reason and a source model the operator didn't form, and a model a year behind the frontier can supply both.
